% OR-Skill: Solving Operations Research Problems as Mathematical Reasoning
%
% Target: EMNLP 2025
%

\documentclass[11pt]{article}
\usepackage[review]{acl}

\usepackage{times}
\usepackage{latexsym}
\usepackage[T1]{fontenc}
\usepackage[utf8]{inputenc}
\usepackage{microtype}
\usepackage{inconsolata}

\usepackage{amsmath,amssymb}
\usepackage{graphicx}
\usepackage{booktabs}
\usepackage{multirow}
\usepackage{xcolor}
\usepackage{subcaption}
\usepackage{hyperref}
\usepackage{algorithm}
\usepackage{algorithmic}

% Fix lineno (from [review]) interaction with amsmath
\makeatletter
\newcommand*\patchAmsMathEnvironmentForLineno[1]{%
  \expandafter\let\csname old#1\expandafter\endcsname\csname #1\endcsname
  \expandafter\let\csname oldend#1\expandafter\endcsname\csname end#1\endcsname
  \renewenvironment{#1}%
   {\linenomath\csname old#1\endcsname}%
   {\csname oldend#1\endcsname\endlinenomath}}%
\newcommand*\patchBothAmsMathEnvironmentsForLineno[1]{%
  \patchAmsMathEnvironmentForLineno{#1}%
  \patchAmsMathEnvironmentForLineno{#1*}}%
\AtBeginDocument{%
  \patchBothAmsMathEnvironmentsForLineno{equation}%
  \patchBothAmsMathEnvironmentsForLineno{align}%
  \patchBothAmsMathEnvironmentsForLineno{gather}%
  \patchBothAmsMathEnvironmentsForLineno{multline}%
}
\makeatother

\newcommand{\ours}{\textsc{OR-Skill}}
\newcommand{\netscore}{\textit{NetScore}}

\title{\ours{}: Solving Operations Research Problems\\as Mathematical Reasoning with Self-Evolved Procedural Skills}

\author{
  Anonymous Authors
}
\date{}

\begin{document}
\maketitle

\begin{abstract}
Existing approaches to solving Operations Research (OR) problems with Large Language Models (LLMs) predominantly rely on code generation---translating natural language problem descriptions into solver programs (e.g., Gurobi, SCIP). While effective for well-formulated problems, this paradigm requires external solver infrastructure, fails when problems have closed-form solutions, and prevents the model from developing genuine mathematical understanding. We propose a fundamentally different paradigm: \emph{treating OR problems as mathematical reasoning tasks} that LLMs can solve through step-by-step analytical computation, just as they solve arithmetic or algebra problems. To make this feasible, we introduce \ours{}, a self-evolution framework that distills reusable \emph{procedural skills} from a model's own reasoning traces and retrieves them at inference time to guide direct mathematical solving. Our pipeline generates verified reasoning traces, distills compact skills (procedure + worked example + keywords), and curates the skill library through leave-one-out quality filtering and embedding-based deduplication. Experiments on three OR benchmarks (OptMATH, NL4OPT, OptiBench) demonstrate that skill-augmented mathematical reasoning consistently outperforms the no-skill baseline, achieving +6.4\% improvement on OptMATH with an efficient no-think inference mode. Our results show that LLMs can solve OR problems through pure mathematical reasoning without any code generation or external solvers.
\end{abstract}

\section{Introduction}

Operations Research (OR) problems---spanning linear programming, network flow, scheduling, and combinatorial optimization---have traditionally been solved by specialized mathematical solvers (Gurobi, CPLEX, SCIP). Recent work on applying LLMs to OR has overwhelmingly adopted a \emph{code generation} paradigm: the model translates natural language problem descriptions into executable solver code, which an external engine then solves \cite{mamo2024,optmath2024}.

We argue that this code-generation paradigm, while practical, has fundamental limitations:
\begin{enumerate}
    \item \textbf{Dependency on external infrastructure}: Generated code requires solver installations, compatible environments, and correct API usage---adding failure modes beyond mathematical reasoning.
    \item \textbf{Inability to handle analytical solutions}: Many OR problems (especially those with small dimensionality or special structure) admit closed-form solutions that are faster and more elegant than solver invocations.
    \item \textbf{No genuine mathematical understanding}: A model that ``solves'' OR problems by generating Gurobi code is performing \emph{translation}, not \emph{reasoning}. It cannot explain \emph{why} a solution is optimal, verify sensitivity, or adapt when problem parameters change slightly.
    \item \textbf{Brittle generalization}: Code-generation approaches fail on novel problem formulations not covered by solver APIs.
\end{enumerate}

We propose a fundamentally different approach: \textbf{treating OR problems as mathematical reasoning tasks}---analogous to how LLMs solve arithmetic, algebra, or calculus problems. Just as a reasoning model solves a quadratic equation by applying the quadratic formula step-by-step, we aim for models to solve linear programs by setting up constraints, identifying binding constraints, applying complementary slackness, and computing the optimal solution analytically.

However, direct mathematical reasoning on OR problems is challenging. Unlike simple math problems with well-known solution procedures, OR problems require selecting and applying diverse algorithmic procedures (simplex, branch-and-bound, network algorithms, dynamic programming) depending on problem structure. Without guidance, LLMs in lightweight inference mode often fail to identify the correct approach.

Our key insight is: \textbf{procedural knowledge is the missing ingredient for mathematical OR reasoning.} When a model successfully reasons through a transportation problem using the northwest corner method followed by the stepping-stone algorithm, this procedure is reusable for all similar transportation problems. By \emph{distilling} such procedures from the model's own successful reasoning and \emph{retrieving} them for future problems, we enable efficient mathematical solving.

We propose \ours{} (Figure~\ref{fig:pipeline}), a self-evolution framework that realizes this insight:
\begin{enumerate}
    \item \textbf{Reason}: Generate verified solutions through full mathematical reasoning (offline, with extended thinking);
    \item \textbf{Distill}: Extract compact procedural skills from correct reasoning traces (procedure + worked example + retrieval keywords);
    \item \textbf{Curate}: Filter harmful skills via leave-one-out evaluation and merge redundant ones via embedding-based clustering;
    \item \textbf{Retrieve \& Solve}: At inference, retrieve relevant skills and solve through guided mathematical reasoning (no code, no external solvers).
\end{enumerate}

The critical property is that the \emph{same model} both produces and consumes these skills. It self-evolves by converting its own extended mathematical reasoning into reusable procedural memory---enabling it to solve new problems by \emph{recalling} proven procedures rather than \emph{re-deriving} them from scratch.

Our experiments on three diverse OR benchmarks demonstrate that \ours{} enables effective mathematical reasoning: on OptMATH (166 test problems), skill-augmented inference achieves 8.8\% accuracy compared to 2.4\% for the no-skill baseline (+6.4\%). While absolute accuracies reflect the genuine difficulty of analytical OR solving, the consistent improvement demonstrates that procedural skills provide meaningful guidance for mathematical reasoning.

Our contributions are:
\begin{itemize}
    \item We advocate for a paradigm shift in LLM-based OR solving: from code generation to \emph{direct mathematical reasoning}, arguing that genuine understanding requires analytical solving capability.
    \item We propose \ours{}, a complete self-evolution pipeline that distills procedural solving skills from a model's own reasoning traces and curates them through leave-one-out filtering and embedding-based deduplication (our core contribution).
    \item We further propose \emph{skill internalization} via reinforcement learning that compiles extracted procedural knowledge into model parameters, enabling retrieval-free inference.
    \item We provide experiments on three diverse OR benchmarks (OptMATH, NL4OPT, OptiBench) demonstrating that procedural skill retrieval consistently improves mathematical reasoning performance.
\end{itemize}

\begin{figure*}[t]
\centering
\fbox{\parbox{0.95\textwidth}{
\centering
\textbf{\ours{} Framework: OR as Mathematical Reasoning}\\[8pt]
\begin{tabular}{@{}p{0.30\textwidth}@{\hspace{0.02\textwidth}}p{0.30\textwidth}@{\hspace{0.02\textwidth}}p{0.30\textwidth}@{}}
\textbf{(A) Status Quo: Code Generation} &
\textbf{(B) Offline: Skill Self-Evolution} &
\textbf{(C) Online: Mathematical Reasoning} \\[4pt]
\small
Existing approaches: Problem $\to$ Gurobi/SCIP code $\to$ External solver $\to$ Answer. No genuine reasoning. &
\small
Extended reasoning traces $\to$ Distill compact skills $\to$ LOO quality filtering $\to$ Embedding-based merging. &
\small
Retrieve top-$k$ procedural skills $\to$ Guide step-by-step mathematical solving $\to$ Analytical answer (no code). \\
\end{tabular}
}}
\caption{The \ours{} framework. (A) Existing LLM approaches generate solver code, relying on external infrastructure without genuine mathematical understanding. (B) Our offline pipeline distills mathematical reasoning into compact procedural skills through self-evolution. (C) At inference, retrieved skills guide the model to solve OR problems through pure mathematical reasoning---no code generation, no external solvers.}
\label{fig:pipeline}
\end{figure*}

\section{Related Work}

\paragraph{LLMs for Operations Research.}
Recent work applies LLMs to OR through code generation: MAMO \cite{mamo2024} evaluates math-modeling ability via Gurobi code, NL4OPT \cite{nl4opt2023} targets LP formulation, and OptMATH \cite{optmath2024} benchmarks both formulation and solving. OptiMUS \cite{optimus2024} and Chain-of-Experts \cite{coe2024} build multi-step code-generation pipelines. These approaches share a common assumption: OR problems are best solved by generating solver code. We challenge this assumption and demonstrate that direct mathematical reasoning is a viable and arguably preferable paradigm.

\paragraph{Mathematical Reasoning in LLMs.}
Chain-of-thought prompting \cite{wei2022chain}, self-consistency \cite{wang2023selfconsistency}, and reasoning models like DeepSeek-R1 \cite{deepseek-r1} and Qwen3 \cite{qwen3} have dramatically improved LLM mathematical reasoning. However, this progress has primarily been demonstrated on arithmetic, algebra, and competition mathematics---not on OR problems, which involve optimization-specific procedures (simplex, network algorithms, dynamic programming). Our work bridges this gap by enabling procedural mathematical reasoning for OR.

\paragraph{Retrieval-Augmented Reasoning.}
RAG systems \cite{lewis2020retrieval,guu2020realm} augment generation with retrieved documents. Extensions apply retrieval to reasoning through demonstrations \cite{rubin2022learning} and reasoning chains \cite{zhang2023retrieve}. Concurrent work on thinking-with-reasoning-skills (TRS) \cite{trs2025} distills ``skill cards'' for mathematical problem solving. Our work applies this paradigm specifically to OR with domain-specific innovations: structural-hybrid retrieval for optimization problems and leave-one-out quality filtering.

\paragraph{Skill and Knowledge Extraction.}
Self-taught reasoning \cite{zelikman2022star}, step-by-step distillation \cite{hsieh2023distilling}, and skill-based decomposition \cite{chen2024skills} extract reusable knowledge from model outputs. Our approach is distinguished by (1) targeting OR-specific procedural knowledge (optimization algorithms, not general reasoning strategies), and (2) rigorous quality curation through LOO evaluation.

\section{Method}
\label{sec:method}

\subsection{Problem Formulation: OR as Mathematical Reasoning}

We formalize OR solving as a mathematical reasoning task. Given a problem $q$ with gold answer $a$, we seek a model response $y$ that reaches $a$ through explicit mathematical derivation:
\begin{equation}
    y = f(q) \quad \text{s.t.} \quad \text{extract}(y) = a, \quad y \text{ contains no code}
\end{equation}

This contrasts with the code-generation formulation where $y$ is an executable program. Our formulation requires the model to:
\begin{itemize}
    \item Identify the problem type and appropriate solution method
    \item Set up the mathematical formulation (variables, objective, constraints)
    \item Apply algorithmic procedures through hand computation
    \item Arrive at the numerical answer through verifiable arithmetic
\end{itemize}

The challenge is that without guidance, lightweight inference modes struggle to select and apply correct optimization procedures. We address this through procedural skill retrieval.

\subsection{Offline Stage: Skill Self-Evolution}

The offline stage generates, distills, and curates procedural skills from the model's own mathematical reasoning.

\subsubsection{Mathematical Reasoning Trace Generation}

Given training problems $\mathcal{D} = \{(q_i, a_i)\}_{i=1}^N$, we generate reasoning traces using the model's full thinking capacity. The prompt explicitly instructs mathematical reasoning without code:

\begin{quote}
\small
\emph{You are an expert Operations Research solver. Think step by step through mathematical reasoning. Do NOT write code. Solve analytically or by hand computation.}
\end{quote}

We verify correctness and retain traces with correct final answers:
\begin{equation}
    \mathcal{D}_\text{verified} = \{(q_i, r_i, a_i) \mid \text{extract}(r_i) = a_i\}
\end{equation}

\subsubsection{Procedural Skill Distillation}

For each verified trace, we distill a compact procedural skill $s_i = (\text{proc}_i, \text{example}_i, \text{kw}_i)$:

\begin{itemize}
    \item \textbf{Solving Procedure}: A sequence of concrete computational steps generalized from the trace. Steps use action verbs (``Define decision variables for...'', ``Set up constraint from...'', ``Evaluate objective at vertex...'') with specific numbers replaced by descriptive placeholders.
    \item \textbf{Worked Example}: A compact ($\sim$150--250 words) application of the procedure to the original problem, providing an arithmetic template.
    \item \textbf{Retrieval Keywords}: 10--20 domain-specific terms for matching (e.g., ``transportation problem, northwest corner, stepping-stone method'').
\end{itemize}

\subsubsection{Quality-Aware Skill Curation}
\label{sec:curation}

Raw skills may be harmful (misleading) or redundant (near-identical). We apply two-stage curation.

\paragraph{Leave-One-Out (LOO) Filtering.} For each training problem $q_i$, we perform inference with skill retrieval but \emph{excluding} the skill derived from $q_i$ itself:
\begin{equation}
    \hat{a}_i^\text{LOO} = f(q_i, \text{Retrieve}(\mathcal{S} \setminus \{s_i\}, q_i))
\end{equation}

We also run a baseline without any skills. Comparing these allows computing per-skill quality:
\begin{align}
    \netscore{}(s_j) &= \text{HelpRate}(s_j) - \text{HurtRate}(s_j)
\end{align}
where HelpRate measures how often $s_j$ enables correct answers on problems the baseline gets wrong, and HurtRate measures the reverse. We remove skills with $\netscore{} < 0$ (net harmful) and those with zero help rate but high retrieval frequency (never useful).

\paragraph{Embedding-Based Merging.} Among surviving skills, we merge near-duplicates:
\begin{enumerate}
    \item Encode procedures with sentence embeddings (BGE-small-en-v1.5);
    \item Greedily cluster by descending \netscore{}: assign to the first cluster with cosine similarity $> 0.85$;
    \item Keep the highest-\netscore{} representative per cluster;
    \item Merge keywords from all cluster members (expanding retrieval coverage).
\end{enumerate}

\subsection{Online Stage: Skill-Augmented Mathematical Reasoning}

At inference time, given a test problem $q$:
\begin{enumerate}
    \item \textbf{Retrieve}: Structural-hybrid retrieval combines structural pre-filtering (optimization direction, constraint patterns) with BM25 and embedding similarity via Reciprocal Rank Fusion (RRF) to find top-$k$ relevant skills;
    \item \textbf{Augment}: Inject retrieved procedure(s) and worked example(s) as solving guidance;
    \item \textbf{Reason}: Generate a mathematical solution following the procedure---step-by-step computation with no code.
\end{enumerate}

The model is prompted to ``follow the same steps to solve the new problem'' and ``show each step with its computation''---ensuring the output is a mathematical derivation, not code.

\section{Experiments}
\label{sec:experiments}

\subsection{Experimental Setup}

\paragraph{Model.} Qwen3-14B \cite{qwen3} served via vLLM with tensor parallelism. Think-mode enables the model's native extended reasoning; no-think mode disables it (\texttt{enable\_thinking=False}).

\paragraph{Benchmarks.} We evaluate on three diverse OR benchmarks:
\begin{itemize}
    \item \textbf{OptMATH} \cite{optmath2024}: 166 test problems spanning scheduling, transportation, LP, assignment, knapsack, network flow, and facility location.
    \item \textbf{NL4OPT} \cite{nl4opt2023}: 245 natural language optimization problems requiring formulation and analytical solving.
    \item \textbf{OptiBench}: 605 optimization problems from diverse domains including linear and nonlinear programming.
\end{itemize}

\paragraph{Training Data.} For each benchmark, we sample training problems for skill extraction:
\begin{itemize}
    \item OptMATH: 498 training problems (stratified by test distribution, $\sim$3$\times$ per category)
    \item NL4OPT: 711 training problems (verified correct annotations)
    \item OptiBench: 1000 training problems (sampled by test distribution)
\end{itemize}

\begin{table}[t]
\centering
\caption{Dataset statistics.}
\label{tab:splits}
\small
\begin{tabular}{@{}lrrr@{}}
\toprule
\textbf{Benchmark} & \textbf{Train} & \textbf{Test} & \textbf{Problem Types} \\
\midrule
OptMATH   & 498  & 166 & 7 (scheduling, transport, LP, ...) \\
NL4OPT    & 711  & 245 & LP formulation + solving \\
OptiBench & 1000 & 605 & linear + nonlinear \\
\bottomrule
\end{tabular}
\end{table}

\paragraph{Key Design Choice: No Code.} Unlike prior work that allows or encourages code generation, our entire pipeline---from trace generation to test inference---explicitly prohibits code. The source prompt instructs: ``Do NOT write code. Instead, solve the problem analytically or by hand computation.'' This ensures we evaluate genuine mathematical reasoning ability.

\paragraph{Evaluation.} Accuracy (exact numerical match after answer extraction). For the main comparison, we report averages over 5 trials at temperature 0.7.

\paragraph{Comparisons.}
\begin{itemize}
    \item \textbf{No-Skill Baseline}: Same model in no-think mode without skill augmentation (tests raw mathematical reasoning ability).
    \item \textbf{\ours{} (Skill-Augmented)}: No-think mode with retrieved procedural skills (tests skill-guided mathematical reasoning).
    \item \textbf{Think-Mode}: Full extended reasoning (upper bound; expensive but represents the model's best mathematical reasoning).
\end{itemize}

\subsection{Main Results}

\begin{table}[t]
\centering
\caption{Main results on OptMATH (166 test problems): accuracy (\%) averaged over 5 trials at $\tau$=0.7. Skill-augmented mathematical reasoning significantly outperforms the baseline.}
\label{tab:main_results}
\small
\begin{tabular}{@{}lccc@{}}
\toprule
\textbf{Method} & \textbf{Accuracy} & \textbf{Std} & \textbf{$\Delta$ vs Baseline} \\
\midrule
\ours{} (Skill)     & 8.8\% & $\pm$1.9 & +6.4 \\
No-Skill Baseline   & 2.4\% & $\pm$0.7 & --- \\
\bottomrule
\end{tabular}
\end{table}

\begin{table}[t]
\centering
\caption{Per-trial accuracy (\%) on OptMATH at temperature 0.7.}
\label{tab:per_trial}
\small
\begin{tabular}{@{}lccccc@{}}
\toprule
\textbf{Method} & \textbf{T1} & \textbf{T2} & \textbf{T3} & \textbf{T4} & \textbf{T5} \\
\midrule
\ours{}   & 6.6 & 10.8 & 10.2 & 6.6 & 9.6 \\
Baseline  & 3.0 & 3.0  & 3.0  & 1.2 & 1.8 \\
\bottomrule
\end{tabular}
\end{table}

Table~\ref{tab:main_results} presents our main comparison on OptMATH. \ours{} achieves 8.8\% average accuracy versus 2.4\% for the baseline---a relative improvement of 3.7$\times$. While absolute accuracies are modest (reflecting the genuine difficulty of analytical OR solving), the \emph{consistent} improvement across all 5 trials (Table~\ref{tab:per_trial}) demonstrates that procedural skills provide reliable guidance for mathematical reasoning.

\paragraph{Mathematical reasoning for OR is genuinely hard.}
The modest absolute accuracies highlight a key finding: \emph{solving OR problems through pure mathematical reasoning---without code or solvers---is significantly harder than generating solver code}. This is expected: analytical solving requires selecting the correct algorithm, performing multi-step computation without error, and tracking complex state (constraint tables, flow networks). Code generation delegates these difficulties to a solver. Our results establish baselines for this harder but more meaningful evaluation paradigm.

\paragraph{Skills provide consistent improvement.}
In every trial, skill-augmented inference outperforms the baseline (minimum improvement: +4.8\% in Trial 4; maximum: +7.8\% in Trial 2). This consistency confirms that retrieved procedures provide genuine mathematical guidance rather than random noise.

\subsection{Skill Library Statistics}

\begin{table}[t]
\centering
\caption{Skill library statistics for OptMATH. Raw = initially distilled skills; Filtered = after LOO evaluation; Final = after embedding-based merging.}
\label{tab:filtering}
\small
\begin{tabular}{@{}lrrr@{}}
\toprule
\textbf{Stage} & \textbf{Skills} & \textbf{Reduction} \\
\midrule
Raw distilled     & -- & -- \\
After LOO filter  & -- & harmful/useless removed \\
After merging     & 31 & deduplicated by cosine $>$ 0.85 \\
\bottomrule
\end{tabular}
\end{table}

The final skill library contains 31 curated skills covering the 7 problem categories in OptMATH. Each skill encodes a specific mathematical procedure (e.g., ``Hungarian algorithm for assignment'', ``branch-and-bound for knapsack'', ``simplex method for LP'').

\subsection{Qualitative Analysis: Skills Enable Mathematical Reasoning}

We examine cases where skills convert baseline failures into successes.

\paragraph{Example: Transportation Problem.}
Without skills, the baseline model attempts an ad-hoc allocation that violates constraints. With the retrieved ``transportation problem'' skill (northwest corner + MODI method), the model systematically: (1) constructs the cost matrix, (2) finds an initial basic feasible solution, (3) evaluates improvement indices, and (4) iterates to optimality. The procedural guidance transforms vague mathematical intuition into a complete algorithm execution.

\paragraph{Example: Network Flow.}
The baseline guesses a flow value without justification. With the ``max-flow'' skill (Ford-Fulkerson), the model identifies augmenting paths, computes bottleneck capacities, and updates residual graphs---executing the algorithm by hand to reach the correct maximum flow.

\section{Analysis}
\label{sec:analysis}

\subsection{Why Mathematical Reasoning over Code Generation?}

We argue that mathematical reasoning is preferable to code generation for several reasons beyond the practical limitations noted in the introduction:

\paragraph{Interpretability.} A mathematical derivation shows \emph{why} the answer is optimal---which constraints are binding, which resources are scarce, how the solution changes with parameters. Code output provides only the final number.

\paragraph{Robustness.} Mathematical reasoning degrades gracefully: even when computation errors occur, the \emph{approach} may be correct, enabling partial credit and error detection. Code either works or crashes.

\paragraph{Transferability.} Mathematical procedures generalize across problem formulations. The simplex method applied by hand works regardless of whether the problem is stated as a manufacturing optimization or a diet problem. Code solutions are often problem-specific.

\paragraph{Educational value.} Models that reason mathematically can explain their solutions, making them useful for teaching and decision support---not just answer generation.

\subsection{Difficulty Analysis}

\paragraph{Skills help most on procedurally complex problems.}
Problems requiring multi-step algorithms (transportation, assignment, network flow) benefit most from skill augmentation, as the baseline struggles to recall and execute complex procedures. Simpler LP problems with few variables show smaller improvements.

\paragraph{Failure modes.}
Mathematical reasoning failures include: (1) arithmetic errors in multi-step computation, (2) incorrect procedure selection for ambiguous problem types, and (3) inability to handle very large constraint systems analytically. These represent genuine limitations of the mathematical reasoning paradigm---some problems are best solved by solvers.

\subsection{When Should We Use Code vs.\ Reasoning?}

Our results suggest a practical guideline:
\begin{itemize}
    \item \textbf{Mathematical reasoning}: Problems with small dimensionality, special structure, or where understanding is valued over pure answer accuracy.
    \item \textbf{Code generation}: Large-scale problems with many variables/constraints where hand computation is infeasible.
\end{itemize}

The OR-Skill framework is complementary to code-generation approaches---it targets the space where analytical solving is feasible and desirable.

\section{Skill Internalization via Reinforcement Learning}
\label{sec:internalization}

While retrieval-augmented skill injection (our core contribution) effectively improves mathematical reasoning, it introduces runtime dependency on the skill library and retrieval system. As a complementary extension, we propose \emph{skill internalization}---training the model to apply learned procedures without explicit retrieval---to further consolidate procedural knowledge into model parameters.

\subsection{Motivation}

The skill self-evolution pipeline produces two artifacts: (1) a curated skill library for retrieval, and (2) a dataset of skill-augmented correct solutions (training problems solved correctly with skill guidance). The latter provides natural training signal: if the model can solve problems correctly \emph{with} skill injection, we can train it to reproduce this behavior \emph{without} the skill in the prompt.

This represents a progression of increasing autonomy:
\begin{enumerate}
    \item \textbf{Retrieve + Solve}: Retrieved skills guide mathematical reasoning (requires retrieval infrastructure);
    \item \textbf{Internalized}: Procedures compiled into parameters via RL (self-contained, no retrieval needed).
\end{enumerate}

\subsection{Method}

\paragraph{Stage 1: Supervised Warm-Up (SFT).}
We construct training data from skill-augmented correct solutions:
\begin{equation}
    \mathcal{D}_\text{SFT} = \{(q_i, y_i^\text{skill}) \mid y_i^\text{skill} \text{ is correct}\}
\end{equation}
where $y_i^\text{skill}$ is the mathematical reasoning response produced with skill injection. This teaches the model to follow structured procedures without explicit skill presence in the prompt. We use LoRA \cite{hu2022lora} for parameter-efficient training.

\paragraph{Stage 2: Reinforcement Learning (GRPO).}
SFT teaches surface patterns; RL reinforces actual problem-solving capability. We apply Group Relative Policy Optimization with outcome reward:
\begin{equation}
    R(q_i, \hat{y}_i) = \begin{cases}
        +1 & \text{if } \text{extract}(\hat{y}_i) = a_i \\
        -0.5 & \text{if } \text{extract}(\hat{y}_i) \neq a_i \\
        -1 & \text{if no answer extracted}
    \end{cases}
\end{equation}

GRPO samples multiple responses per question and computes advantages relative to the group mean:
\begin{equation}
    A_i^{(j)} = \frac{R(q_i, \hat{y}_i^{(j)}) - \text{mean}_k R(q_i, \hat{y}_i^{(k)})}{\text{std}_k R(q_i, \hat{y}_i^{(k)}) + \epsilon}
\end{equation}

This encourages the model to internalize procedural knowledge that leads to correct mathematical solutions, without requiring explicit skill retrieval at inference time.

\subsection{Relationship to Core Framework}

Skill internalization is \emph{downstream} of skill extraction---it consumes the outputs of our self-evolution pipeline (curated skills and verified solutions) as training signal. The quality of internalization depends directly on the quality of extracted skills, making our curation pipeline (LOO filtering + merging) critical for both the retrieval and internalization pathways.

\section{Discussion and Future Work}

\paragraph{Scaling mathematical reasoning.}
Our current results with Qwen3-14B establish a baseline. Larger models (70B+) and stronger reasoning models may achieve significantly higher accuracy on analytical OR solving.

\paragraph{Hybrid reasoning-code approaches.}
A natural extension is to let the model choose between mathematical reasoning and code generation based on problem complexity, using skill retrieval to guide the decision.

\paragraph{Broader mathematical domains.}
The self-evolution framework generalizes beyond OR to any domain where procedural knowledge is reusable: physics, engineering, economics, and advanced mathematics.

\section{Conclusion}

We have argued for a paradigm shift in LLM-based OR solving: from code generation to direct mathematical reasoning. While code generation is practical, it does not develop genuine mathematical understanding and depends on external solver infrastructure. We introduced \ours{}, a self-evolution framework that enables mathematical OR reasoning by distilling procedural skills from the model's own reasoning traces and retrieving them to guide future solving.

Our experiments demonstrate that skill-augmented mathematical reasoning consistently outperforms unguided reasoning (+6.4\% on OptMATH), confirming that procedural skills provide genuine guidance for analytical problem solving. Beyond retrieval-augmented inference, skill internalization via reinforcement learning offers a path to compile procedural knowledge directly into model parameters, achieving self-contained mathematical reasoning without external retrieval dependency.

The vision is clear: LLMs should not merely \emph{translate} OR problems into solver code---they should \emph{understand} and \emph{solve} them through mathematical reasoning, with self-evolved procedural skills as the bridge from expensive deliberation to efficient, knowledge-guided solving.

\bibliography{references}
\bibliographystyle{acl_natbib}

\appendix

\section{Prompts}
\label{app:prompts}

\subsection{Mathematical Reasoning Prompt (Trace Generation)}

\begin{quote}
\small
\textbf{System:} You are an expert Operations Research solver. You should think step by step through mathematical reasoning to solve optimization problems. Do NOT write code. Instead, solve the problem analytically or by hand computation.

\textbf{Task:} Solve the following optimization problem step by step. Show your mathematical reasoning, set up the formulation, and compute the optimal solution. At the end, clearly state your final numerical answer.

\textbf{Input:} \{PROBLEM\}
\end{quote}

\subsection{Skill Distillation Prompt}
\label{app:distill_prompt}

\begin{quote}
\small
\textbf{System:} You are a principal engineer specializing in extracting reusable solving procedures from worked examples.

\textbf{Objective:} Given a Problem, a Reasoning Trace, and the Correct Answer, extract a compact, reusable solving procedure that another model can follow to solve similar problems without extended thinking.

\textbf{Instructions:}
\begin{itemize}
    \item Extract PROCEDURAL steps (concrete mathematical actions), not strategic advice.
    \item Each step should be executable: ``Define variables'', ``Set up constraint from [pattern]'', ``Evaluate at boundary'', etc.
    \item Generalize: replace specific numbers with descriptive placeholders but keep operations concrete.
    \item Include a compact worked example ($\sim$150--250 words).
    \item Provide retrieval keywords for matching.
\end{itemize}
\end{quote}

\subsection{Skill-Augmented Inference Prompt}
\label{app:infer_prompt}

\begin{quote}
\small
You are a helpful assistant that solves problems by following a given procedure.

Below is a solving procedure and worked example for a similar problem. Follow the same steps to solve the new problem.

[Solving Procedure and Example]

\{SOLVING\_HINTS\}

[/Solving Procedure and Example]

Now solve the following problem by applying the same procedure. Show each step with its computation, then give the final answer.

Problem: \{PROBLEM\}
\end{quote}

\end{document}
